\documentclass[a4paper,12pt]{article}
\usepackage[utf8]{inputenc}
\usepackage[spanish]{babel}

\title{Informe municipal de seguimiento de incidencias}
\author{Area de Atencion Ciudadana}
\date{\today}

\begin{document}

\maketitle

\section{Objetivo}

Este informe ficticio sirve para practicar la creacion de documentos estructurados con LaTeX en Visual Studio Code.

El documento resume incidencias comunicadas por la ciudadania sobre elementos de la via publica. Los datos son inventados y se usan solo con fines formativos.

\section{Resumen de incidencias}

Durante la semana analizada se han registrado incidencias en tres ambitos principales:

\begin{itemize}
    \item Alumbrado publico: 8 avisos.
    \item Aceras y pavimento: 5 avisos.
    \item Limpieza viaria: 6 avisos.
\end{itemize}

\section{Actuaciones previstas}

Las actuaciones se organizaran en funcion de la urgencia, la disponibilidad de medios y la concentracion de avisos en una misma zona.

\begin{itemize}
    \item Revisar los avisos de alumbrado en calles con mayor numero de incidencias.
    \item Priorizar reparaciones de pavimento que afecten a la accesibilidad peatonal.
    \item Coordinar la limpieza de zonas con avisos repetidos.
\end{itemize}

\section{Comprobaciones antes de enviar}

Antes de usar un informe de este tipo en un contexto real, una persona debe comprobar:

\begin{itemize}
    \item que los datos proceden de una fuente autorizada;
    \item que no se incluyen datos personales innecesarios;
    \item que las cifras son correctas;
    \item que las actuaciones propuestas han sido validadas por el area responsable.
\end{itemize}

\end{document}
